%======================================================================
%===  DTU Poster — SHRED Robustness Study
%======================================================================
\documentclass[
     title     = {{Robust Field Reconstruction from Sparse Sensors: SHRED vs SDN vs QR-POD under Sensor Corruption}}
    ,author    = {{Kasper van der Horst and Dimitris Papantzikos}}
    ,subject   = {{Data-Driven Methods for Engineers --- Spring 2026}}
    ,bgcolor   = dtucoolgrey
    ,highlight = dtured
    ,toplogo   = {{tex_dtu_logo_white}}
    ,botlogo   = {{tex_dtu_compute_b_uk}}
    ,papersize = {{a0paper}}
    ,colcount  = {{3columns}}
    ,longtitle
    ,largecaption
]{dtuposter}

\usepackage[T1]{fontenc}
\usepackage[utf8x]{inputenc}
\usepackage{cmbright}
\usepackage{arevmath}
\renewcommand{\familydefault}{\sfdefault}
\usepackage{enumitem}
\setlist{nosep,leftmargin=*,topsep=2pt}
\usepackage{booktabs}
\usepackage{multirow}
\usepackage{graphicx}
\usepackage{amsmath}
\usepackage{adjustbox}
\usepackage{caption}

% Reduce title to ~65pt (longtitle default is 100pt)
% This \AtBeginDocument runs after the class's own hook, overriding the komafont
\makeatletter
\AtBeginDocument{%
  \setkomafont{title}{\normalfont\fontsize{65}{78}\selectfont\bfseries\color{\dtuposter@highlight}}%
  \setkomafont{minisec}{\normalfont\fontsize{42}{38}\selectfont\bfseries}%
  \setkomafont{dictum}{\normalfont\fontsize{34}{32}\selectfont}%
}
\renewenvironment{dtuposterhead}{%
  \flushleft {\usekomafont{title} \dtuposter@title } \par \vspace*{0.75\baselineskip}%
}{%
  \vspace*{0.75\baselineskip} \par%
}
\makeatother


\begin{document}

%======================================================================
%===  Header
\begin{dtuposterhead}
\dtuposterauthor{Kasper van der Horst \& Dimitris Papantzikos}
\dtuposteraffil{Technical University of Denmark, DTU Compute, Kgs.\ Lyngby, Denmark}
\end{dtuposterhead}

%======================================================================
\begin{dtupostercontent}

%==============================================================
% COLUMN 1 — background & methods
%==============================================================

\section{Motivation}

High-dimensional physical fields (e.g.\ fluid vorticity) are expensive
to measure densely. A few strategically placed sensors can drive a
learned decoder to reconstruct the full field---provided the decoder
tolerates real-world sensor faults such as noise, signal loss, and
burst interference. Which in theory should be possible as SHRED performs well
at low number of sensors, motivitating the idea that even when not all sensors
provide meaningfull readings, a good reconstruction should still be feasbible.

We benchmark three reconstruction methods across sensor budgets and
corruption types, and assess whether corruption-matched augmentation
during training improves robustness without hurting clean accuracy.

\section{Methods}

\textbf{SHRED} (SHallow REcurrent Decoder) [1] feeds a window of
$L{=}10$ sensor readings into a two-layer LSTM (hidden size 64),
then a three-layer MLP ($350{\to}400{\to}m$) predicts all $m$ field
values simultaneously. Temporal context confers implicit robustness:
corruption of a single timestep has limited effect on the hidden state.

\textbf{SDN} ablates temporal history: one sensor snapshot feeds the
same MLP directly, isolating the LSTM's contribution.

\textbf{QR-POD} is a linear gappy-POD baseline. Sensor locations are
chosen by QR-pivoted POD column selection. Reconstruction is an
analytic least-squares solve.

\section{Augmentation Strategies}

Applied to sensor inputs during \emph{training only}; validation and
test always use clean inputs:
\begin{itemize}
  \item \emph{Clean}: no augmentation (baseline)
  \item \emph{Gaussian}: additive noise ($\sigma{=}0.03$) on a random subset
  \item \emph{Dropout}: random sensors zeroed out (one always survives)
  \item \emph{Hybrid}: noise, zero-out, or nothing with equal probability
  \item \emph{Burst}: temporally correlated burst-noise on random sensors
\end{itemize}

\section{Experimental Setup}

2-D cylinder wake DNS at $Re{=}100$: $N{=}5{,}000$ snapshots on a
$128{\times}64$ grid. Sequential split (train 4\,970 / val 20 / test 10).
Sensor placement via QR-pivoted POD ($r{=}20$ retained modes).
Adam optimiser, MSE loss, $lr{=}10^{-3}$, batch 64, early stopping
with patience $5{\times}20{=}100$ validation epochs, best-weight restore.

%==============================================================
% COLUMN 2 — figure + references
%==============================================================

\section{Robustness Comparison (QR Placement)}

\noindent\includegraphics[width=\linewidth]{%
  ../../cylinder/robustness/robustness_bar_QR}
\captionof{figure}{Relative $\ell_2$ error for QR sensor placement across all
  corruption scenarios and sensor budgets (3, 5, 8 sensors).
  \textbf{Solid}: clean test set.
  \textbf{Hatched}: corrupted test set.
  Each group pairs the clean-trained model with its corruption
  specialist for SHRED, SDN, and QR-POD.}
\label{fig:robbar}

%==============================================================
% COLUMN 3 — key findings + table + conclusions
%==============================================================

\section{Key Findings}

\begin{itemize}
  \item SHRED and SDN beat QR-POD on \textbf{clean data} at every sensor
        budget---by more than $10\times$ at 3 sensors (SHRED: 0.006 vs
        QR-POD: 0.227).
  \item Under \textbf{Gaussian noise}, clean-trained SHRED degrades only
        mildly (0.006${\to}$0.010 at 5 sensors), suggesting the LSTM
        averages out noise implicitly.
  \item \textbf{Dropout} is the hardest corruption: clean models collapse
        to $0.18$--$0.42$ error; specialist models recover to
        $0.032$--$0.107$.
  \item \textbf{Hybrid and Burst} are well handled by specialist SHRED
        (errors $0.030$--$0.093$) with minimal clean-data penalty.
  \item \textbf{QR-POD robustness} is competitive under Gaussian noise at
        8 sensors (0.065), but is dominated everywhere else by the
        neural methods.
  \item Overall: while using data augmentation in training does show improvements. The results are still below what is achieved using SHRED and clean training data.
\end{itemize}

\section{Burst Reconstruction --- 5 Sensors, QR Placement}

\noindent\includegraphics[width=\linewidth]{%
  ../../cylinder/robustness/sensors_5_QR/panel_burst}
\captionof{figure}{Field reconstruction under burst sensor corruption (5 sensors, QR placement).
  Rows show ground truth, clean-trained SHRED, burst-specialist SHRED, SDN, and QR-POD.}
\label{fig:burst}

\section{The limitations of SHRED}

While Shred is shown to perform very well on clean data, even for very low number of sensors. There is a key limitation in the way the SHRED architecture handles especially sensor dropout.
During training it has to learn specific mappings from the inputs of sensors to the output through the LSTM and SDN. The current architecture of shred is not able to effectively learn conficting sensor dropouts which have to map to the same output.
So over the full training it learns a compromise between these readings, a compromise which is suboptimal for each seperate configuration.
The only way for SHRED to be able to handle this is for the input to the LSTM to be a better representation of the physical system than the raw sensor inputs. This however is a highly nontrivial issue. 
We did implement a variational state space approach extension to SHRED where the input was mapped to a latent space using an small encoder network.
This approach proved highly unstable in training yielding greatly varying results across different seeds. Thus, this approach cannot yet be claimed to provide a solution to this problem. 
\section{References}

[1]~J.\ Williams, O.\ Zahn, J.\ N.\ Kutz, ``Sensing with shallow
recurrent decoder networks,'' \textit{Proc.\ R.\ Soc.\ A}
480:20230816, 2024.

\end{dtupostercontent}

\end{document}
